\documentclass[11pt,a4paper]{article}

\usepackage[utf8]{inputenc}
\usepackage[T1]{fontenc}
\usepackage{lmodern}
\usepackage{microtype}
\usepackage[french]{babel}
\usepackage[margin=1.8cm,top=2.1cm,bottom=2.1cm]{geometry}
\usepackage{xcolor}
\usepackage{amsmath,amssymb}
\usepackage{booktabs}
\usepackage{tabularx}
\usepackage{eurosym}
\usepackage{tcolorbox}
\tcbuselibrary{skins,breakable}
\usepackage{tikz}
\usetikzlibrary{shapes.geometric,arrows.meta,positioning,calc,fit,backgrounds}
\usepackage{listings}
\usepackage{hyperref}
\usepackage{fancyhdr}
\usepackage{titlesec}
\usepackage{enumitem}

% Palette de couleurs professionnelle
\definecolor{primary}{RGB}{20, 65, 105}       % Bleu marine
\definecolor{secondary}{RGB}{35, 100, 140}    % Bleu acier
\definecolor{accent}{RGB}{215, 110, 10}       % Ambre marketing
\definecolor{darkgray}{RGB}{40, 50, 65}
\definecolor{lightgray}{RGB}{246, 248, 251}
\definecolor{bordergray}{RGB}{220, 226, 235}
\definecolor{success}{RGB}{20, 95, 50}
\definecolor{successbg}{RGB}{240, 253, 244}
\definecolor{warningbg}{RGB}{254, 243, 199}
\definecolor{warningframe}{RGB}{180, 83, 9}

\hypersetup{
    colorlinks=true,
    linkcolor=primary,
    urlcolor=secondary,
    citecolor=accent,
    pdftitle={Guide de Conception - Agent Metier Marketing SMA},
    pdfauthor={Antigravity}
}

% En-tête et pied de page
\pagestyle{fancy}
\fancyhf{}
\setlength{\headheight}{14pt}
\rhead{\textcolor{darkgray}{\footnotesize \textbf{Guide SMA} -- Agent Métier Marketing}}
\lhead{\textcolor{darkgray}{\footnotesize Ingénierie des Systèmes Multi-Agents}}
\rfoot{\textcolor{darkgray}{\footnotesize Page \thepage}}
\renewcommand{\headrulewidth}{0.4pt}
\renewcommand{\footrulewidth}{0.4pt}

% Espacement titres
\titlespacing*{\section}{0pt}{1.1ex plus 0.3ex minus 0.2ex}{0.5ex plus 0.2ex}
\titlespacing*{\subsection}{0pt}{0.8ex plus 0.2ex minus 0.2ex}{0.4ex plus 0.2ex}
\titlespacing*{\subsubsection}{0pt}{0.6ex plus 0.2ex minus 0.2ex}{0.3ex plus 0.2ex}

\titleformat{\section}
  {\normalfont\large\bfseries\color{primary}}{\thesection.}{0.7em}{}[\titlerule]
\titleformat{\subsection}
  {\normalfont\normalsize\bfseries\color{secondary}}{\thesubsection}{0.5em}{}
\titleformat{\subsubsection}
  {\normalfont\small\bfseries\color{darkgray}}{\thesubsubsection}{0.4em}{}

% Boîtes d'alerte et concepts
\newtcolorbox{conceptbox}[1]{
  colback=lightgray,
  colframe=primary,
  fonttitle=\bfseries\small\color{white},
  title={#1},
  arc=2mm,
  boxrule=0.8pt,
  left=3mm, right=3mm, top=2mm, bottom=2mm,
  before skip=5pt, after skip=5pt
}

\newtcolorbox{alertbox}[1]{
  colback=successbg,
  colframe=success,
  fonttitle=\bfseries\small\color{white},
  title={#1},
  arc=2mm,
  boxrule=0.8pt,
  left=3mm, right=3mm, top=2mm, bottom=2mm,
  before skip=5pt, after skip=5pt
}

\newtcolorbox{warningbox}[1]{
  colback=warningbg,
  colframe=warningframe,
  fonttitle=\bfseries\small\color{white},
  title={#1},
  arc=2mm,
  boxrule=0.8pt,
  left=3mm, right=3mm, top=2mm, bottom=2mm,
  before skip=5pt, after skip=5pt
}

\newtcolorbox{marketingbox}[1]{
  colback=orange!5!white,
  colframe=accent,
  fonttitle=\bfseries\small\color{white},
  title={#1},
  arc=2mm,
  boxrule=0.8pt,
  left=3mm, right=3mm, top=2mm, bottom=2mm,
  before skip=5pt, after skip=5pt
}

% Configuration Listings (Code)
\lstset{
  backgroundcolor=\color{lightgray},
  basicstyle=\ttfamily\scriptsize\color{darkgray},
  keywordstyle=\bfseries\color{primary},
  commentstyle=\itshape\color{gray},
  stringstyle=\color{accent},
  breaklines=true,
  frame=single,
  rulecolor=\color{bordergray},
  tabsize=2,
  showstringspaces=false,
  numbers=left,
  numberstyle=\tiny\color{gray},
  numbersep=6pt,
  aboveskip=4pt,
  belowskip=4pt
}

\begin{document}

% ==================== PAGE 1 ====================
\begin{center}
    {\LARGE \textbf{\color{primary}Guide de Conception d'un Agent Métier Marketing}}\\[0.15cm]
    {\large \textbf{\color{secondary}Fondements SMA, Négociation FIPA-ACL \& Intégration au Noyau Central}}\\[0.25cm]
    \rule{0.85\textwidth}{0.8pt}\\[0.15cm]
    {\footnotesize \textbf{Document de Cadrage Méthodologique et Technique pour l'Équipe Projet} | Semestre Académique}
\end{center}
\vspace{-0.25cm}

\begin{conceptbox}{Objectif Global du Projet}
Votre équipe a pour mission de développer un \textbf{Agent Logiciel de Marketing}. Dans la simulation, votre agent est responsable de valoriser, promouvoir et échanger des \textbf{produits finis} avec les agents conçus par les autres équipes (acheteurs, usines de fabrication ou distributeurs).

Pour orchestrer ce marché, votre enseignant fournit un \textbf{Noyau central} (middleware / bus d'échange). Ce document synthétise tout ce que vous devez impérativement maîtriser pour débuter sereinement : concepts fondamentaux, fonctionnement du noyau, stratégie marketing, code d'intégration et pièges à éviter.
\end{conceptbox}

\section{Comprendre les Agents et les Systèmes Multi-Agents (SMA)}

\subsection{Qu'est-ce qu'un Agent logiciel ?}
Contrairement à un script classique ou à une API passive qui ne s'exécute que lorsqu'on l'appelle explicitement, un \textbf{Agent Intelligent} est une entité logicielle autonome située dans un environnement dynamique, dotée de buts propres et capable d'agir de son propre chef.

\begin{center}
\small
\begin{tabularx}{\textwidth}{p{3cm} X X}
\toprule
\textbf{Critère} & \textbf{Programme Classique / Script} & \textbf{Agent Intelligent (SMA)} \\
\midrule
\textbf{Contrôle} & Passif : exécute l'ordre reçu sans délibérer. & Autonome : décide \textit{si} et \textit{comment} il répond. \\
\textbf{Déclenchement} & Appel de méthode synchrone externe. & Boucle autonome perpétuelle d'écoute et d'initiative. \\
\textbf{Environnement} & Entrées prévisibles et contrôlées. & Marché ouvert, concurrentiel et incertain. \\
\textbf{Comportement} & Déterministe (\textit{Entrée} $\rightarrow$ \textit{Sortie}). & Stratégique (\textit{Perception} $\rightarrow$ \textit{Délibération} $\rightarrow$ \textit{Action}). \\
\bottomrule
\end{tabularx}
\end{center}

\subsection{Les 4 Piliers d'un Agent (Modèle de Wooldridge \& Jennings)}
\begin{itemize}[leftmargin=*, itemsep=1.5pt]
    \item \textbf{Autonomie :} L'agent contrôle son état interne, ses fonds et ses choix sans dépendre d'un humain.
    \item \textbf{Réactivité :} L'agent perçoit son environnement (offres reçues, alertes de marché) et réagit à temps.
    \item \textbf{Proactivité :} L'agent ne subit pas ; il prend l'initiative de démarcher des clients et de diffuser des offres.
    \item \textbf{Sociabilité :} L'agent communique et négocie via un protocole formel et structuré (actes de parole).
\end{itemize}

\subsection{La Boucle Fondamentale : Perception \texorpdfstring{$\rightarrow$}{->} Délibération \texorpdfstring{$\rightarrow$}{->} Action}
Le cœur d'exécution de l'agent tourne en continu selon trois étapes fondamentales :

\begin{center}
\begin{tikzpicture}[
    box/.style={draw=primary, fill=lightgray, rounded corners=1.5mm, font=\bfseries\footnotesize, align=center, inner sep=4pt, minimum height=0.85cm},
    arrow/.style={->, >=Stealth, thick, color=secondary}
]
    \node[box, fill=primary!15] (percept) at (0, 0) {1. PERCEPTION\\ \scriptsize Réception message / Marché};
    \node[box, fill=accent!20] (delib) at (5.2, 0) {2. DÉLIBÉRATION\\ \scriptsize Stratégie Marketing / Prix};
    \node[box, fill=success!20] (action) at (10.4, 0) {3. ACTION\\ \scriptsize Envoi PROPOSE, ACCEPT, CFP};
    \node[box, fill=darkgray!15, minimum width=11.6cm] (env) at (5.2, -1.35) {ENVIRONNEMENT : Noyau du Professeur + Marché des Autres Équipes};

    \draw[arrow] (env.north -| percept.south) -- (percept.south) node[midway, left, font=\scriptsize] {Événements};
    \draw[arrow] (percept.east) -- (delib.west);
    \draw[arrow] (delib.east) -- (action.west);
    \draw[arrow] (action.south) -- (env.north -| action.south) node[midway, right, font=\scriptsize] {Messages / Réponses};
\end{tikzpicture}
\end{center}

\subsection{Le Modèle BDI (Beliefs, Desires, Intentions)}
Pour structurer la mémoire et le raisonnement de l'agent, le modèle d'architecture de référence est le modèle \textbf{BDI} :
\begin{itemize}[leftmargin=*, itemsep=1.5pt]
    \item \textbf{Beliefs (Croyances / Connaissances) :} Ce que l'agent a en mémoire à l'instant $t$ (stock disponible, trésorerie, historique des prix des concurrents, réputation des partenaires).
    \item \textbf{Desires (Désirs / Buts) :} Les finalités économiques (ex: écouler le stock au cours de la simulation en maximisant la marge nette).
    \item \textbf{Intentions (Plans engagés) :} L'action précise en cours d'exécution (ex: négocier une concession de 5\% sur le lot \#4).
\end{itemize}

% ==================== PAGE 2 ====================
\section{Le Rôle du Noyau (« Kernel ») fourni par le Professeur}

Le noyau fourni par votre enseignant ne concourt pas comme une équipe ; c'est \textbf{l'infrastructure d'arbitrage et de transport (Middleware)} du Système Multi-Agents.

\subsection{Les 4 Fonctions Primordiales du Noyau}
\begin{enumerate}[leftmargin=*, itemsep=2pt]
    \item \textbf{Annuaire et Découverte (Directory Facilitator / Pages Jaunes) :} Registre officiel des agents connectés. Votre agent interroge l'annuaire pour découvrir qui vend et qui achète.
    \item \textbf{Bus de Messagerie et Routage :} Chaque agent se connecte \textbf{uniquement au noyau}. Le noyau achemine les messages de façon centralisée et sécurisée.
    \item \textbf{Arbitrage et Validation des Transactions (Registre Officiel) :} Si un contrat de vente est signé, le noyau valide que le vendeur possède les produits et que l'acheteur a les crédits nécessaires.
    \item \textbf{Gestion du Temps (Clock / Rounds) :} La simulation progresse par tours synchronisés (\textit{rounds} discrets cadencés) ou en flux continu asynchrone.
\end{enumerate}

\subsection{Le Protocole de Dialogue : Standard FIPA-ACL}
Les échanges respectent la grammaire standard FIPA-ACL (\textit{Agent Communication Language}), formulée en \textbf{JSON} :

\begin{center}
\small
\begin{tabularx}{\textwidth}{p{3.8cm} p{4.2cm} X}
\toprule
\textbf{Performative (Acte)} & \textbf{Intention du Message} & \textbf{Exemple d'Usage Marketing} \\
\midrule
\texttt{CFP} (\textit{Call For Proposal}) & Appel d'offres / Sollicitation & « Avis au marché : qui souhaite acquérir 20 unités de Produit-A ? » \\
\texttt{PROPOSE} & Soumission d'une offre chiffrée & « Je propose 20 unités au tarif unitaire de 45\,\euro. » \\
\texttt{ACCEPT\_PROPOSAL} & Accord et conclusion du contrat & « Offre validée, nous confirmons la transaction. » \\
\texttt{REJECT\_PROPOSAL} & Déclin de la proposition & « Prix inacceptable, proposition déclinée. » \\
\texttt{INFORM} & Transmission d'informations & « Notification : mise à jour de notre catalogue commercial. » \\
\bottomrule
\end{tabularx}
\end{center}

\subsection{Protocole Type de Négociation : Contract Net Protocol (CNP)}
Voici le schéma de négociation bilatérale le plus fréquent lors de la vente de produits finis :

\begin{center}
\begin{tikzpicture}[
    node distance=1.8cm, auto,
    actor/.style={draw=primary, fill=primary!10, rounded corners=1.5mm, font=\bfseries\footnotesize, minimum width=2.6cm, minimum height=0.65cm},
    msg/.style={->, >=Stealth, thick, font=\scriptsize}
]
    \node[actor] (mkt) at (0, 0) {Votre Agent (Marketing)};
    \node[actor] (noyau) at (5.2, 0) {Noyau (Middleware)};
    \node[actor] (client) at (10.4, 0) {Agent Client (Autre Équipe)};

    \draw[dashed, gray] (mkt) -- (0, -3.1);
    \draw[dashed, gray] (noyau) -- (5.2, -3.1);
    \draw[dashed, gray] (client) -- (10.4, -3.1);

    \draw[msg, color=primary] (0, -0.6) -- (5.2, -0.6) node[midway, above] {1. \texttt{CFP} (Diffusion du lot)};
    \draw[msg, color=primary] (5.2, -0.8) -- (10.4, -0.8) node[midway, above] {Transmission};

    \draw[msg, color=accent] (10.4, -1.3) -- (5.2, -1.3) node[midway, above] {2. \texttt{PROPOSE} (Offre d'achat: 38\,\euro)};
    \draw[msg, color=accent] (5.2, -1.5) -- (0, -1.5) node[midway, above] {Transmission};

    \draw[msg, color=success] (0, -2.0) -- (5.2, -2.0) node[midway, above] {3. \texttt{ACCEPT} ou contre-offre \texttt{PROPOSE}};
    \draw[msg, color=success] (5.2, -2.2) -- (10.4, -2.2) node[midway, above] {Transmission};

    \draw[msg, color=secondary] (5.2, -2.8) -- (0, -2.8) node[midway, above] {4. \texttt{INFORM} (Confirmation transaction par le Noyau)};
\end{tikzpicture}
\end{center}

\begin{conceptbox}{Format Concret d'une Trame JSON FIPA-ACL}
Tout message transitant par le noyau possède une structure standardisée :
\begin{lstlisting}[language=Python]
{
  "sender": "Equipe_Marketing_3",
  "receiver": "Equipe_Acheteur_7",
  "conversation_id": "conv_lot_942",
  "performative": "PROPOSE",
  "content": {"product": "Module_Camera_HD", "quantity": 10, "unit_price": 42.50}
}
\end{lstlisting}
\end{conceptbox}

% ==================== PAGE 3 ====================
\section{L'Agent Métier Marketing : Stratégies et Algorithmes Décisionnels}

\subsection{Missions Spécifiques de l'Agent Marketing}
Dans la simulation, votre agent incarne la \textbf{direction commerciale} de votre équipe :
\begin{itemize}[leftmargin=*, itemsep=2pt]
    \item \textbf{Positionnement Produit :} Calculer la valeur marchande des produits finis selon leurs caractéristiques.
    \item \textbf{Pricing Dynamique :} Ajuster les prix selon le coût de revient, la demande globale et l'échéance du temps.
    \item \textbf{Tactique de Concession :} Savoir accorder des rabais calculés pour conclure sans brader la rentabilité.
    \item \textbf{Veille Concurrentielle :} Consigner l'historique des transactions pour adapter ses marges aux autres équipes.
\end{itemize}

\subsection{Modèle Mathématique de Tarification Dynamique Temporelle}
Soit $C$ le coût unitaire de fabrication, $P_{\max}$ le prix d'ouverture désiré, et $P_{\min}$ le prix plancher en dessous duquel l'équipe vend à perte ($P_{\min} > C$). Au tour ou temps $t \in [0, T]$, le prix cible $P(t)$ est calculé par :
\begin{equation}
P(t) = P_{\min} + (P_{\max} - P_{\min}) \cdot \left(1 - \frac{t}{T}\right)^{\beta}
\end{equation}
\begin{itemize}[leftmargin=*, itemsep=2pt]
    \item \textbf{$\beta > 1$ (Stratégie Concessionnaire) :} Baisse rapide des prix dès le début pour liquider les stocks.
    \item \textbf{$\beta = 1$ (Stratégie Linéaire) :} Concession régulière et constante au fil des tours.
    \item \textbf{$0 < \beta < 1$ (Stratégie Tenace / Écrémage) :} Maintien de prix élevés le plus longtemps possible.
\end{itemize}

\subsection{Matrice Décisionnelle Face à une Offre Reçue}
Quand un agent client vous soumet une offre de prix $P_{\text{offre}}$, la délibération applique cette règle formelle :
\begin{equation}
\text{Décision}(P_{\text{offre}}) = 
\begin{cases} 
\texttt{ACCEPT\_PROPOSAL} & \text{si } P_{\text{offre}} \ge P(t) \\
\texttt{PROPOSE}\left( \frac{P(t) + P_{\text{offre}}}{2} \right) & \text{si } P_{\min} \le P_{\text{offre}} < P(t) \quad \text{(Contre-offre de compromis)} \\
\texttt{REJECT\_PROPOSAL} & \text{si } P_{\text{offre}} < P_{\min}
\end{cases}
\end{equation}

\section{Architecture Logicielle \& Guide d'Intégration au Noyau}

\subsection{Architecture Modulaire en 4 Couches}
Pour rester sereins et agiles face aux évolutions du code du professeur, découpez votre agent en 4 couches hermétiques :

\begin{center}
\begin{tikzpicture}[
    layer/.style={draw=primary, fill=lightgray, rounded corners=1.5mm, text width=12.5cm, font=\scriptsize, minimum height=0.68cm, align=left, inner sep=4pt},
    arrow/.style={<->, >=Stealth, thick, color=accent}
]
    \node[layer, fill=accent!15] (l4) at (0, 2.4) {\textbf{Couche 4 : Stratégie Métier Marketing} (Pricing dynamique, matrice de décision, concession)};
    \node[layer, fill=secondary!15] (l3) at (0, 1.6) {\textbf{Couche 3 : État \& Mémoire (Beliefs)} (Portefeuille, stock de produits finis, réputation)};
    \node[layer, fill=primary!15] (l2) at (0, 0.8) {\textbf{Couche 2 : Protocole FIPA-ACL} (Désérialisation JSON $\leftrightarrow$ Objets métiers)};
    \node[layer, fill=darkgray!15] (l1) at (0, 0.0) {\textbf{Couche 1 : Connecteur Réseau} (Client WebSockets, Socket TCP ou API HTTP vers le Noyau)};

    \draw[arrow] (l1.north) -- (l2.south);
    \draw[arrow] (l2.north) -- (l3.south);
    \draw[arrow] (l3.north) -- (l4.south);
\end{tikzpicture}
\end{center}

\subsection{Checklist Technique : Les questions à poser impérativement au professeur}
\begin{alertbox}{Checklist d'Intégration Réseau}
\begin{itemize}[leftmargin=*, itemsep=1.5pt]
    \item[\textbf{[ ]}] \textbf{Moyen de transport :} WebSockets, Sockets TCP brutes, API REST HTTP, bus MQTT ou wrapper Python fourni ?
    \item[\textbf{[ ]}] \textbf{Schéma des données :} Clés JSON exactes (\texttt{"performative"}, \texttt{"sender"}, \texttt{"receiver"}, \texttt{"content"}) ?
    \item[\textbf{[ ]}] \textbf{Handshake :} Quel message exact envoyer pour s'enregistrer auprès de l'annuaire au lancement ?
    \item[\textbf{[ ]}] \textbf{Cadencement temporel :} Simulation par tours (\textit{rounds}) discrets ou temps réel asynchrone ?
    \item[\textbf{[ ]}] \textbf{Structure d'un produit :} Quels attributs définissent un produit final (ID, qualité, coût unitaire) ?
\end{itemize}
\end{alertbox}

% ==================== PAGE 4 ====================
\section{Implémentation de Référence (Template Python Modulaire)}

Ce squelette Python complet illustre la séparation absolue entre la logique de décision marketing et la tuyauterie réseau :

\begin{lstlisting}[language=Python]
import json

class MarketingStrategy:
    """Couche 4 (Metier) : Pricing dynamique et tactique de negociation."""
    def __init__(self, cost_price, min_margin=1.15, max_margin=1.70):
        self.cost = cost_price
        self.p_min = cost_price * min_margin      # Prix plancher absolu (rentabilite)
        self.p_max = cost_price * max_margin      # Prix d'ouverture ideal (forte marge)

    def compute_target_price(self, progress_ratio, beta=1.5):
        """Concession temporelle selon l'avancement du round (0.0 a 1.0)."""
        return self.p_min + (self.p_max - self.p_min) * ((1.0 - progress_ratio) ** beta)

    def evaluate_offer(self, offered_price, progress_ratio):
        target = self.compute_target_price(progress_ratio)
        if offered_price >= target:
            return "ACCEPT_PROPOSAL", offered_price
        elif offered_price >= self.p_min:
            counter_offer = round((target + offered_price) / 2.0, 2)
            return "PROPOSE", counter_offer
        return "REJECT_PROPOSAL", None

class MarketingAgent:
    """Couches 2 & 3 : Gestion de l'etat interne et reaction aux messages."""
    def __init__(self, agent_id, team_id, connector):
        self.agent_id = agent_id
        self.team_id = team_id
        self.connector = connector
        self.strategy = MarketingStrategy(cost_price=50.0)
        self.inventory = {}
        self.negotiation_history = []

    def register_to_kernel(self):
        """Enregistrement officiel aupres du noyau."""
        payload = {"performative": "REGISTER", "sender": self.agent_id, "team": self.team_id}
        self.connector.send(json.dumps(payload))

    def on_message_received(self, raw_json):
        """Boucle fondamentale : Perception -> Deliberation -> Action"""
        msg = json.loads(raw_json)
        performative = msg.get("performative")
        sender = msg.get("sender")
        content = msg.get("content", {})

        if performative == "CFP":
            # Reponse proactive a un appel d'offres entrant
            opening_price = self.strategy.p_max
            reply = {
                "sender": self.agent_id,
                "receiver": sender,
                "performative": "PROPOSE",
                "content": {"product_id": content.get("product_id"), "price": opening_price}
            }
            self.connector.send(json.dumps(reply))

        elif performative == "PROPOSE":
            # Evaluation d'une proposition de prix recue
            decision, price = self.strategy.evaluate_offer(content.get("price", 0), progress_ratio=0.5)
            reply = {
                "sender": self.agent_id,
                "receiver": sender,
                "performative": decision,
                "content": {"price": price} if price else {}
            }
            self.connector.send(json.dumps(reply))

        elif performative == "INFORM":
            # Confirmation d'une transaction executee par le noyau
            print(f"[{self.agent_id}] Transaction confirmee par le noyau : {content}")
\end{lstlisting}

\begin{conceptbox}{Comment Exécuter ce Code dès Aujourd'hui ?}
Vous n'avez pas besoin d'attendre le serveur du professeur pour tester ! Il suffit d'écrire une fausse classe \texttt{connector} avec une méthode \texttt{send(msg)} qui affiche le message à l'écran, et d'appeler manuellement \texttt{on\_message\_received(json\_test)}.
\end{conceptbox}

% ==================== PAGE 5 ====================
\section{Feuille de Route pour l'Équipe Projet}

Pour avancer méthodiquement sans vous bloquer mutuellement, appliquez ce découpage en 4 phases :

\begin{enumerate}[leftmargin=*, itemsep=2.5pt]
    \item \textbf{Sprint 1 : L'Agent Squelette (Dummy Agent)}
    \begin{itemize}[itemsep=1pt]
        \item Mettre en place le dépôt Git et l'architecture en 4 couches.
        \item Écrire le connecteur réseau minimal : connexion au noyau, envoi du message \texttt{REGISTER}, et confirmation des messages reçus (« ping/pong »).
    \end{itemize}

    \item \textbf{Sprint 2 : Le Simulateur Factice (« Mock Agent » local)}
    \begin{itemize}[itemsep=1pt]
        \item Coder un petit script Python simulant un agent client acheteur sur votre machine locale.
        \item \textbf{Avantage clé :} Permet à l'équipe de coder et tester la négociation sans attendre la disponibilité du noyau du professeur !
    \end{itemize}

    \item \textbf{Sprint 3 : Implémentation du Cœur Marketing}
    \begin{itemize}[itemsep=1pt]
        \item Développer la fonction de prix dynamique et la gestion des concessions.
        \item Mémoriser l'historique des prix des autres équipes pour optimiser les marges.
    \end{itemize}

    \item \textbf{Sprint 4 : Intégration Inter-Équipes \& Résilience}
    \begin{itemize}[itemsep=1pt]
        \item Déploiement grandeur nature avec les autres équipes sur le réseau universitaire.
        \item Gestion robuste des anomalies (timeouts, messages mal formés, déconnexions).
    \end{itemize}
\end{enumerate}

\begin{marketingbox}{Conseil Pratique : Répartition des Rôles au sein de l'Équipe}
\begin{itemize}[leftmargin=*, itemsep=1.5pt]
    \item \textbf{Étudiant 1 (Lead Réseau \& Noyau) :} Connecteur (Sockets/WebSockets), sérialisation JSON, boucle d'écoute.
    \item \textbf{Étudiant 2 (Lead Métier Marketing) :} Algorithmes de pricing, matrice de décision, stratégies de concession.
    \item \textbf{Étudiant 3 (Lead QA / Simulation) :} Script de simulation factice (Mock), tests unitaires, validation des échanges.
\end{itemize}
\end{marketingbox}

\section{Les 5 Pièges Fréquents en Projet SMA \& Leurs Solutions}

\begin{warningbox}{Attention aux Erreurs Classiques d'Étudiants}
\begin{enumerate}[leftmargin=*, itemsep=2pt]
    \item \textbf{Le Blocage Synchrone (Deadlock) :} Ne faites jamais un \texttt{socket.recv()} sans timeout au milieu de votre logique métier. Si un agent adverse plante, votre agent sera gelé. \textit{Solution : Utilisez des écoutes asynchrones ou des threads séparés.}
    \item \textbf{La Boucle Infinie de Négociation (Ping-Pong) :} Deux agents qui se renvoient sans cesse des contre-offres sans progresser font planter la simulation. \textit{Solution : Ajoutez un compteur d'échanges (ex: max 3 tours de négociation par transaction).}
    \item \textbf{La Vente à Perte par Concession Excessive :} Si vous concédez trop vite sous la panique du temps, vous vendrez en dessous de votre coût de revient $C$. \textit{Solution : Imposez une barrière matérielle inviolable dans le code : \texttt{assert price >= p\_min}.}
    \item \textbf{Le Double Engagement de Stock :} Proposer le même lot de produits à deux acheteurs simultanés conduit à un échec d'arbitrage. \textit{Solution : Réservez temporairement le stock lors de l'émission d'une proposition (\texttt{reserved\_stock}).}
    \item \textbf{Attendre le Noyau du Professeur :} L'erreur la plus coûteuse est de rester passif. \textit{Solution : Développez immédiatement votre stratégie marketing en local avec des tests unitaires.}
\end{enumerate}
\end{warningbox}

\end{document}
